\newpage
\pagenumbering{roman}
\setcounter{page}{2}

\renewcommand{\abstractname}{\large Апстракт}
\phantomsection
%\addcontentsline{toc}{section}{\abstractname}


\begin{abstract}
\fontsize{12pt}{15pt}\selectfont
    \textbf{Прецизната тридимензионална реконструкција е еден од клучните и најистрајните предизвици во полето на компјутерската визија и компјутерската графика. Целта на реконструкцијата е креирање на веродостојна дигитална репрезентација на објекти или сцени од реалниот свет, користејќи податоци добиени од дводимензионални медиуми, како фотографии или видео-записи. Процесот се базира на инверзија на начинот на кој функционира перцепцијата, односно трансформација на дводимензионални податоци во тридимензионални структури разбирливи за компјутерските системи.
    \\
    Во рамките на оваа димпломска работа се анализираат трите водечки пристапи на реконструкција: класичната техника на фотограметрија, како и современите пристапи базирани на машинско учење: Neural Radiance Fields (NeRF) и Gaussian splatting. Ќе се испитуваат нивните методолошки пристапи, потреби од податоци, пресметковна ефикасност и визуелен квалитет на конечните резултати, со цел да се добие подлабоко разбирање на предностите и недостатоците на поединечните методи за најпрецизно да се одлучи кој е најсоодветен метод за користење во различни сценарија. Посебен акцент се става на нивната примена во здравството, археологијата, дигитализацијата на културното наследство, виртуелната и проширената реалност, како и потенцијалната комерцијална употреба.}

    
\fontsize{10pt}{12pt}\selectfont
    \textbf{Клучни зборови:} \textit{Фотограметрија, Neural radiance fields, Gaussian splatting, тридимензионална реконструкција, машинско учење, дигитализација}
\end{abstract}
